\SourceLeaf{415}%
nach \([6]\) die folgende Formel anwenden
\[\BareRoot{\left\{rr-\left(\frac{g\mp F}{\lambda}\right)^2\right\}}\]
\par
Auch hier liefert eine Wurzel also zwei Lösungen, nämlich die symmetrischen Lagen der Punkte \((1),(2),(3)\) auf entgegengesetzten Seiten von \(\mathrm{(II)}\) in gleichem Abstand. Für \(x=+1\) hat die positive Richtung auf \(\mathrm{(I)}\) denselben Sinn wie auf \(\mathrm{(II)}\), für \(x=-1\) den entgegengesetzten. Ausgenommen werden muss nur der Fall, in dem ohne Rücksicht auf das Vorzeichen \(\lambda r=g-F\) (für \(x=+1\)) oder \(=g+F\) (für \(x=-1\)) gilt. Dann fallen die beiden symmetrischen Lagen von \(\mathrm{(I)}\) zusammen und stimmen mit der Geraden \(\mathrm{(II)}\) selbst überein.

Von den Wurzeln der Gleichung \([7]\) sind offenbar die imaginären und die reellen außerhalb der Grenzen \(-1\) und \(+1\) auszuschließen, ebenso die Wurzel \(+1\) oder \(-1\) selbst, wenn \(\lambda r\) ohne Rücksicht auf das Vorzeichen kleiner ist als \(g-F\) beziehungsweise \(g+F\). Es lässt sich außerdem beweisen, dass stets mindestens eine der drei Wurzeln zu den auszuschließenden gehört und dass es daher insgesamt nie mehr als vier verschiedene Lösungen mit reellen Koordinaten geben kann. Genau genommen bildet \textit{ein ganz singulärer} Fall insofern eine Ausnahme von der ersten Aussage, als dabei keine Wurzel ausgeschlossen wird. Es ist der bereits erwähnte Fall, in dem für \(x=\pm1\) die Ordinate~\(=0\) wird und in dem, wie sich leicht beweisen lässt, die betreffende Wurzel doppelt zählt. Der Ausdruck links vom Gleichheitszeichen in Gleichung \([7]\) enthält also den Faktor \((x\mp1)^2\). Die Gleichung hat dann nur zwei verschiedene Wurzeln, von denen die zweite durchaus ebenfalls zulässig sein kann. Dieser Ausnahmefall berührt die Schlussfolgerung selbst somit nicht.

Schließlich ist noch zu bemerken, dass es auch unter den Lösungen in reellen Zahlen physisch unzulässige geben kann. Vom Kreis mit Mittelpunkt \((1)\) durch \((4)\) und \((5)\) bildet nur ein bestimmter Bogen den geometrischen Ort des Punktes \((1)\), nämlich der Bogen auf der positiven Seite von \((4),(5)\), wenn \(\vartheta\) kleiner als \(180^0\) ist, und den Bogen auf der negativen Seite, wenn \(\vartheta\) ein überstumpfer Winkel ist. Dasselbe gilt für die beiden anderen Kreise. Diese physische Bedingung ist in unserer Lösung noch nicht berücksichtigt. Von den verschiedenen in reellen Zahlen gefundenen Lösungen sind also nur diejenigen zulässig, bei denen die Werte, die sich für die Abszissen der drei Punkte \((1),(2),(3)\)
